\chapter{Primewise Derivation of the Final Coherence Factorization}

\section{Sequential normalization}
Let
\[
 p^e\parallel c.
\]
The local output contact packet has order
\[
 p^{e-1}.
\]
The final factorization separates this depth into distinct mechanisms. The separation must be sequential because the same prime-power depth can otherwise be counted more than once.

Start with the signed input exponent vector
\[
 \nu=(\nu_q)_{q\mid ab},
 \qquad
 \nu_q=v_q(b)-v_q(a).
\]
Let $K_p$ be the $p$-primary part of the common Kummer factor. Divide the exponent vector and the local modulus by this common power. The normalized vector is primitive relative to the remaining modulus.

For odd $p$, form the logarithmic row
\[
 \ell_{p,e}(q)
 =
 \frac1p\log_p\langle q\rangle_p
 \pmod{p^{e-1}}.
\]
Let
\[
 w_p
 =
 \min_{q\mid ab}
 v_p(\ell_{p,e}(q))
\]
after Kummer normalization. Divide the row and the modulus by $p^{w_p}$. The remaining row is primitive modulo the residual power $p^{\chi_p}$.

The local residual character is
\[
 \chi_p(x)
 =
 \exp\left(
 \frac{2\pi i}{p^{\chi_p}}g_p\cdot x
 \right).
\]
Its kernel contains the normalized exponent vector. The local factors are therefore
\[
 K_p,
 \qquad
 W_p=p^{w_p},
 \qquad
 C_p=p^{\chi_p}.
\]
At two, one sign layer $\epsilon_{2,p}$ is removed before the weighted logarithmic row is formed.

Multiplying over output primes gives
\[
 \boxed{
 M=\epsilon_2KWC.
 }
\]
The equality follows because the local exponents partition the original contact depth:
\[
 e-1
 =
 v_p(\epsilon_2)
 +
 v_p(K)
 +
 v_p(W)
 +
 v_p(C).
\]

\section{Why the Kummer layer comes first}
Suppose every input exponent is divisible by $p^k$. Then the balanced rational packet is already a $p^k$th power in the valuation lattice. The corresponding output contact depth should not also be attributed to a local logarithm defect.

Dividing the exponent vector first removes this preload. The logarithmic row is then evaluated on a primitive packet. Its common coefficient depth measures a genuine degeneration of the logarithm operator rather than a common power already present in the input.

This is the same normalization principle used in regulator theory: separate the algebraic divisibility of the vector before measuring the determinant of the linear form.

\section{Odd-prime coefficient depth}
Let $q\neq p$ be prime. Write
\[
 q=\omega(q)\langle q\rangle_p,
\]
where $\omega(q)$ is the Teichm\"uller factor. Since
\[
 \langle q\rangle_p^{p-1}=q^{p-1},
\]
we have
\[
 \log_p\langle q\rangle_p
 =
 \frac1{p-1}\log_p(q^{p-1}).
\]
If
\[
 s=v_p(q^{p-1}-1)>0,
\]
then
\[
 v_p(\log_p(q^{p-1}))=s.
\]
Therefore
\[
 v_p\left(\frac1p\log_p\langle q\rangle_p\right)
 =s-1.
\]
Truncating at the available contact depth gives
\[
 v_p(\ell_{p,e}(q))
 =
 \min\{e-1,s-1\}.
\]

\section{Dyadic coefficient depth}
For odd $q$, separate the sign by writing
\[
 q=(-1)^{(q-1)/2}\langle q\rangle_2,
 \qquad
 \langle q\rangle_2\in1+4\Z_2.
\]
Then
\[
 \langle q\rangle_2^2=q^2.
\]
The weighted coefficient is
\[
 \ell_{2,e}(q)
 =
 \frac14\log_2\langle q\rangle_2.
\]
If
\[
 s=v_2(q^2-1),
\]
then
\[
 v_2(\ell_{2,e}(q))=s-3,
\]
truncated at $e-2$. The missing first layer is the separate sign bit. This is the final form of the weighted dyadic branch.

\section{CRT synthesis of the global row}
The local residual moduli $C_p$ are pairwise coprime. Let
\[
 C=\prod_{p\mid c}C_p.
\]
For each input coordinate $q$, the Chinese remainder theorem gives a unique residue
\[
 g_q\pmod C
\]
whose reduction modulo $C_p$ is the local coefficient $g_{p,q}$.

The resulting row
\[
 g=(g_q)_{q\mid ab}
\]
is primitive modulo $C$ because every local row is primitive. Its character
\[
 \chi_{\mathrm{coh}}(x)
 =
 e^{2\pi i g\cdot x/C}
\]
has exact order $C$. The coherence lattice is its kernel.

\section{A determinant viewpoint}
When several independent exponent packets are available, replace the one-row character by an exterior-power regulator. Let $L$ be a normalized cross-prime logarithm matrix. A nonzero $r\times r$ minor
\[
 \det L_{I,J}
\]
measures the failure of $r$ input packets to cancel simultaneously across $r$ output channels.

The one-row coherence quotient is the rank-one case. A higher-rank determinant theorem may be stronger because it can use several relations or several independent row constraints at once.

\chapter{Exact Smith and CRT Calculations}

\section{One-row Smith normal form}
For an integer row
\[
 a=(a_1,\ldots,a_s),
\]
let
\[
 d=\gcd(a_1,\ldots,a_s).
\]
There exists a unimodular matrix $U\in\mathrm{GL}_s(\Z)$ such that
\[
 aU=(d,0,\ldots,0).
\]
Therefore the image of the row is $d\Z$, the kernel has rank $s-1$, and the quotient by the kernel is infinite cyclic scaled by $d$.

This elementary fact drives sector compression, local lift gauges, and the primitive-gradient formal-smoothness theorem.

\section{Two-row Smith data}
For a two-row matrix $M$, the first Smith factor is the gcd of all entries. The product of the two factors is the gcd of all nonzero $2\times2$ minors. This allows the obstruction torsion to be computed without constructing the full unimodular reduction.

For the exponent-free phase obstruction matrix, every nonzero cross-sector minor is
\[
 \pm\frac{abc}{pq}
 =
 \pm\frac{abc}{R}\frac{R}{pq}.
\]
All such minors are divisible by $abc/R$. For each prime $\ell\mid R$, choose one column with $p=\ell$ and one column from a different nonempty sector. The reduced minor $R/(pq)$ is then not divisible by $\ell$. Hence the gcd of the reduced minors is one and the second Smith factor is exactly $abc/R$.

\section{Local CRT lift in closed form}
Let
\[
 r=\prod_{i=1}^{m}p_i,
 \qquad
 w_i=r/p_i.
\]
The equation
\[
 \sum_iw_ix_i=t
\]
has a solution because
\[
 \gcd(w_1,\ldots,w_m)=1.
\]
Modulo $p_i$,
\[
 x_i\equiv t w_i^{-1}\pmod{p_i}.
\]
Let $\xi_i$ be the centered residue. Then
\[
 \sum_iw_i\xi_i-t
\]
is divisible by $r$. Set
\[
 q=\frac{t-\sum_iw_i\xi_i}{r}.
\]
Writing $x_i=\xi_i+p_i z_i$ gives
\[
 \sum_i z_i=q.
\]

For a box $|x_i|\le H$, the permitted $z_i$ form intervals. The sum of independent integer intervals is every integer between the sum of lower endpoints and the sum of upper endpoints. This proves the exact feasibility test.

\section{Binary search certificate}
The lower and upper endpoint sums are monotone in $H$. Therefore the least feasible height is found by binary search. The result includes both:

\begin{enumerate}[leftmargin=1.5em]
\item a witness at height $H$
\item an interval-separation certificate at height $H-1$
\end{enumerate}

This is stronger than an optimization routine returning a candidate. It certifies global minimality for the local lift.

\section{Polyphase circulant determinant}
For the valuation word
\[
 V_e(T)=1+T+\cdots+T^{e-1}
\]
in $A_N=\Z[T]/(T^N-1)$, multiplication is circulant. Its determinant is
\[
 \prod_{\zeta^N=1}V_e(\zeta).
\]
Let $I_N$ be the augmentation ideal of $A_N$. When $\gcd(e,N)=1$, the nontrivial roots are permuted by $\zeta\mapsto\zeta^e$, so multiplication by $V_e$ has determinant one on the integral lattice $I_N$ and is therefore an automorphism of it. On $A_N/I_N\cong\Z$, the map is multiplication by $e$. The exact sequence for $I_N\subset A_N$ then gives cokernel $\Z/e\Z$ and Smith form $\operatorname{diag}(1,\ldots,1,e)$.

\section{Synthesis-normalized refinement}
Suppose a coarse integer $t$ is split into $e$ shell coordinates. The minimum child sup norm is roughly $|t|/e$. The shell-sum synthesis map has norm $e$. The product of analysis norm and synthesis norm remains at least one. Any claimed gain must therefore be stated in a synthesis-normalized norm.

\section{Exact ratio semantics}
Three ratios recur in the computation:
\[
 \frac{\Phi}{\overline H},
 \qquad
 \frac{\Phi}{k_*\overline H},
 \qquad
 \frac{H(1)}{\Phi}.
\]
They are all at least one by definition. A value greater than one is a discrete overhead, not an error. Exact rational strings should accompany decimal displays when the integers are large.

\chapter{Finite Data and Interpretation Rules}

\section{What the finite data record}
A useful data row contains at least:

\begin{longtable}{p{0.24\textwidth}p{0.65\textwidth}}
\toprule
Field & Meaning\\
\midrule
$a,b,c$ & primitive additive triple\\
prime factorizations & exact valuation data for all three terms\\
$R$ & global radical\\
$q$ & standard quality\\
$M=c/R_c$ & output contact modulus\\
$\Theta,\Phi$ & coarse and phase packet heights when computed\\
$H(1),k_*,H(k_*)$ & exact moving-search data\\
$\epsilon_2,K,W,C$ & sequential final factorization\\
$\nu_0$ & normalized input exponent vector\\
local rows & normalized odd and dyadic logarithmic filters\\
coherence row & CRT-combined primitive character\\
certificates & witnesses, empty-box bounds, and exact residual checks\\
\bottomrule
\end{longtable}

\section{What a high-quality row does not prove}
A large $C$, a short relation vector, or a striking phase resonance in one row does not imply a uniform law. The row can suggest a mechanism. A theorem requires a family definition and a proof that the mechanism persists or is rare.

\section{Recommended finite experiments}
The most informative next experiments are not simply larger bounds on $c$. They should stratify by mechanism.

\begin{enumerate}[leftmargin=1.5em]
\item Fix the number of input primes and vary output depth.
\item Fix the output prime and search for large local coefficient defects.
\item Separate $W>1$ rows from $W=1$ coherence rows.
\item Measure normalized logarithmic minors rather than only one coherence quotient.
\item Search families generated by recurrences or perfect powers.
\item Compare exact lattice invariants before and after gluing two relations.
\end{enumerate}

\section{Diagnostic interpretation}
Every computational test should state what conclusion it does and does not support.

\begin{longtable}{p{0.28\textwidth}p{0.61\textwidth}}
\toprule
Failure type & Diagnostic\\
\midrule
Kummer preload & common gcd of modulus and valuation vector\\
entrywise logarithmic defect & large generalized Wieferich coefficient depth\\
collective coherence & large residual quotient with ordinary coefficients\\
rank collapse & small logarithmic minors\\
archimedean cost & large period, degree, or synthesis norm\\
integral cost & large denominator, covolume, or congruence exponent\\
out-of-scope object & contact depth not represented in the chosen global carrier\\
\bottomrule
\end{longtable}

\section{From finite observation to theorem}
A finite observation should pass through four mathematical levels.

\begin{enumerate}[leftmargin=1.5em]
\item Numerical observation on a declared finite range.
\item Exact symbolic identity on each tested row.
\item Structural conjecture stated independently of the data.
\item Proof or a precise counterexample family.
\end{enumerate}

This protocol prevents a finite pattern from being described as an established theorem.

\chapter{Theorem Summary and Structural Limitations}

\section{Exact theorem summary}
The exact core can be compressed to the following list.

\begin{enumerate}[leftmargin=1.5em]
\item Prime-packet Leibniz rule.
\item One-row additive lattice.
\item Arithmetic Wronskian and radical divisibility.
\item Three-sector Smith compression.
\item Polyphase Smith form with one aggregate invariant.
\item Joint obstruction torsion $abc/R$.
\item Exact lattice extension $0\to L_0\to L\to\Z\to0$.
\item Exact real stable height.
\item Formal smoothness and all-order integral jet extension.
\item Exact moving-interval search and finite stopping.
\item Prime-shell energy identities.
\item Divisor-polyphase perfect reconstruction.
\item Atkin--Lehner scaled Walsh bank.
\item Integral projector denominator theorem.
\item Filtered contact quotient and carry correction.
\item Exact cyclic gcd transfer.
\item Witt realization of contact depth.
\item Exact Kummer kernel and logarithmic bound.
\item Fermat height compression and character bank.
\item Full-plane Fermat line-bundle cost.
\item Polylogarithmic distribution and three-point closure.
\item Exact finite-place logarithmic notch.
\item Exact Kummer--Wieferich--coherence factorization.
\item Coherence lattice index and cyclic quotient.
\item Arbitrary-index limitation theorem.
\item Direct positive divisibility lower bound.
\item Mean-preserving Hecke cascade.
\end{enumerate}

\section{Conditional implications}
The strongest conditional routes are:

\begin{enumerate}[leftmargin=1.5em]
\item small transverse packet implies $abc$
\item filtered cusp-to-Hecke transfer with two size bounds implies $abc$
\item sparse Kummer--Jacobi quotient with two size bounds implies $abc$
\item coherence conductor bound implies $abc$
\item suitable adelic logarithmic minor bound implies $abc$
\end{enumerate}

The final two are the most direct endpoints.

\section{Structural limitations}
The following routes are closed without a new ingredient:

\begin{enumerate}[leftmargin=1.5em]
\item more exact search of the same first-order invariant
\item more local lift existence
\item unitary shell transforms
\item raw coordinate splitting
\item direct split persistence into one cyclic quotient
\item the full Fermat character bank
\item small complex periods without integral control
\item abstract lattice reduction without arithmetic row structure
\item direct positive contact relations
\item linear Hecke averaging
\end{enumerate}

\section{Arithmetic endpoint}
The arithmetic endpoint is the coherence character
\[
 \chi_{\mathrm{coh}}:
 \Z^{S_{ab}}\to\Z/C\Z
\]
whose row is built from simultaneous logarithms of the actual input primes and whose kernel contains the normalized exponent vector forced by $a+b=c$.

The missing theorem is not generic. It must exploit the shared arithmetic origin of all local rows and the exact identity of the output primes as divisors of the sum.
